\chapter{Design and Methodology}
\label{ch:method}

\section{System Overview}
\label{sec:overview}

\viki\ is an offline pipeline: demonstrations are recorded, processed and
exported in distinct stages, with no requirement that any stage run at
interactive rates. This is a deliberate design choice. Relaxing the real-time
constraint that shapes teleoperation systems permits global optimisation over
whole trajectories, iterative refinement, and a validation stage that executes
on physical hardware---none of which is available to an online system.

The pipeline comprises six stages, shown in Fig.~\ref{fig:pipeline}.

\begin{enumerate}[leftmargin=2em]
  \item \emph{Capture.} Synchronised RGB-D streams are recorded from a
        calibrated multi-camera rig while a demonstrator performs the task
        bare-handed.
  \item \emph{Perception.} Hand landmarks are detected per view, lifted to
        three dimensions using sensor depth, and fused across views with
        confidence weights to yield a wrist pose in $\SE$ and a gripper state.
  \item \emph{Representation.} The wrist trajectory is expressed both relative
        to the manipulated object and relative to the ChArUco workspace anchor.
  \item \emph{Retargeting.} A robot joint trajectory is obtained by minimising
        a single weighted cost functional subject to kinematic and collision
        constraints.
  \item \emph{Replay validation.} The synthesised trajectory is executed on the
        physical manipulator; proprioception is recorded and infeasible
        episodes are rejected.
  \item \emph{Export.} Surviving episodes are written in a LeRobot-compatible
        format for downstream policy training.
\end{enumerate}

\begin{figure}[htbp]
\centering
\begin{tikzpicture}[
  node distance=4mm,
  box/.style={draw, rounded corners=1pt, align=center, font=\footnotesize,
              minimum height=8mm, text width=22mm, inner sep=2pt},
  arr/.style={-{Latex[length=2mm]}, thick}]
\node[box] (cap) {Capture\\\scriptsize 2$\times$Kinect + RealSense};
\node[box, right=of cap] (per) {Perception\\\scriptsize landmarks + fusion};
\node[box, right=of per] (rep) {Representation\\\scriptsize object-relative $\SE$};
\node[box, below=8mm of rep] (ret) {Retargeting\\\scriptsize single cost functional};
\node[box, left=of ret] (val) {Replay\\\scriptsize physical validation};
\node[box, left=of val] (exp) {Export\\\scriptsize LeRobot / HDF5};
\draw[arr] (cap) -- (per);
\draw[arr] (per) -- (rep);
\draw[arr] (rep) -- (ret);
\draw[arr] (ret) -- (val);
\draw[arr] (val) -- (exp);
\draw[arr] (val.north) -- ++(0,4mm) -| node[pos=0.25, above, font=\scriptsize]
  {reject / re-solve} (ret.east);
\end{tikzpicture}
\caption{Stages of the \viki\ pipeline. Capture, perception and representation
are embodiment-agnostic; retargeting, replay validation and export are
specific to the target manipulator. The feedback edge indicates that episodes
failing replay are either re-solved with modified weights or discarded.}
\label{fig:pipeline}
\end{figure}

\section{Hardware and Workspace}
\label{sec:hardware}

The capture rig comprises two Azure Kinect DK cameras and one Intel RealSense
D435i observing a tabletop workspace, with a UR3 manipulator mounted at a fixed
position relative to that workspace. The Azure Kinect was selected for the
primary views because its time-of-flight depth sensing provides substantially
better range estimation than neural regression from RGB, and because its
specified error characteristics---random error standard deviation at most
$17$~mm and systematic error below $11$~mm plus $0.1$ per cent of distance
\cite{mskinect}---are documented and independently verified \cite{tolgyessy}.
The RealSense provides a supplementary close-range view at higher depth
accuracy \cite{intel}.

The two Kinects are placed to preserve substantial overlap in their fields of
view. This is a direct consequence of the finding of Bu and Lee \cite{bulee}
that iterative-closest-point refinement fails between near-opposing views for
lack of point-cloud overlap; near-orthogonal placement retains both
triangulation baseline and overlap.

\section{Calibration}
\label{sec:calib}

\subsection{Intrinsics and extrinsics}

Intrinsic parameters for every colour and depth stream are estimated from
ChArUco board observations. Extrinsic calibration follows the workspace-anchor
principle: rather than expressing cameras relative to one another, each camera
is calibrated against a ChArUco board rigidly fixed to the workspace, defining
a common world frame $W$. For camera $k$ this yields
\begin{equation}
  \bT^{W}_{C_k} \in \SE,
  \label{eq:extrinsic}
\end{equation}
the transform from the camera frame $C_k$ to the workspace frame $W$. Anchoring
to the workspace rather than to a reference camera has two consequences: adding
or moving a camera does not invalidate the calibration of the others, and the
robot base can be registered into the same frame without a separate
camera-to-camera chain.

\subsection{Robot registration}

The manipulator base frame $R$ is registered to $W$ by observing a ChArUco
target held in a known end-effector pose across a set of configurations, giving
$\bT^{W}_{R}$ directly. Where an explicit camera-to-robot transform is required
instead, the Daniilidis dual-quaternion solver is used, following the
robustness evidence of \cite{alihandeye}, with fifteen to twenty poses spanning
a hemisphere.

\section{Synchronisation}
\label{sec:sync}

The two Azure Kinects are hardware-synchronised in master/subordinate
configuration, with subordinate devices started before the master and depth
exposures staggered to avoid mutual laser interference \cite{mskinect}. The
RealSense cannot share that trigger; its synchronisation mode is incompatible
at the electrical level.

Cross-family alignment therefore uses a single host-side monotonic clock. Every
frame, from every device, is stamped on arrival with
\begin{equation}
  \tau = \texttt{time.monotonic\_ns()},
  \label{eq:hostclock}
\end{equation}
and streams are aligned by nearest neighbour in $\tau$. Two points deserve
emphasis. First, per-device hardware timestamps are \emph{not} used for
cross-device arithmetic: distinct device families count their own cycles from
their own epochs, and the vendor documentation states explicitly that equal
timestamps across cameras are not guaranteed \cite{mskinect}. Second, a
monotonic clock rather than wall-clock time is required, since wall-clock time
may be adjusted discontinuously by time synchronisation daemons, which would
corrupt interval computations mid-recording.

Residual alignment jitter is measured rather than assumed
(Section~\ref{sec:res-calib}), and hardware synchronisation of the Kinect pair
is verified by confirming bounded drift in the inter-camera timestamp
difference \cite{ridgerun}.

\section{Perception and Confidence-Weighted Fusion}
\label{sec:fusion}

\subsection{Landmark detection and lifting}

Hand landmarks are detected in each colour stream, yielding for view $k$ at
time $t$ a set of image-plane landmarks with associated visibility scores
$v^{k}_{i,t} \in [0,1]$. Each landmark is lifted to three dimensions using the
registered depth map rather than a learned depth estimate, exploiting the
time-of-flight measurement directly, and transformed into the workspace frame
by \eqref{eq:extrinsic}.

\subsection{Fusion}

Estimates from multiple views are combined by confidence-weighted averaging
rather than by naive mean. For landmark $i$ at time $t$, the fused position is
\begin{equation}
  \hat{\bx}_{i,t} \;=\;
  \frac{\sum_{k} w^{k}_{i,t}\, \bx^{k}_{i,t}}
       {\sum_{k} w^{k}_{i,t}},
  \label{eq:fusion}
\end{equation}
with weights
\begin{equation}
  w^{k}_{i,t} \;=\;
  \underbrace{v^{k}_{i,t}}_{\text{detector}}
  \cdot
  \underbrace{c^{k}_{i,t}}_{\text{sensor}}
  \cdot
  \underbrace{\frac{1}{\left(d^{k}_{i,t}\right)^{2}}}_{\text{range}}
  \cdot
  \underbrace{\max\!\left(0, \cos\theta^{k}_{i,t}\right)}_{\text{incidence}},
  \label{eq:weights}
\end{equation}
where

\begin{center}
\begin{tabular}{@{\hspace{10mm}}p{16mm}@{}p{95mm}}
$v^{k}_{i,t}$ & detector visibility score for landmark $i$ in view $k$;\\
$c^{k}_{i,t}$ & sensor depth confidence or validity flag at that pixel;\\
$d^{k}_{i,t}$ & range from camera $k$ to the landmark;\\
$\theta^{k}_{i,t}$ & angle between the local surface normal and the viewing direction.
\end{tabular}
\end{center}

\noindent The range and incidence factors follow established practice in
geometric depth fusion, where surfaces observed near perpendicular are weighted
above obliquely observed ones because they are measured more accurately
\cite{fusionpatent}; the detector and sensor factors follow confidence-map
re-weighting \cite{zeng}. The product form is chosen so that a zero in any
factor---an undetected landmark, an invalid depth pixel, a grazing view---
removes that view's contribution entirely rather than merely attenuating it.

\subsection{Wrist pose and gripper state}

The wrist pose $\bT^{W}_{H,t} \in \SE$ is obtained by fitting an orthonormal
frame to the fused wrist and knuckle landmarks. The gripper state is derived
from the thumb--index distance,
\begin{equation}
  g_t \;=\;
  \mathbb{1}\!\left[\,\norm{\hat{\bx}_{\text{thumb},t}
  - \hat{\bx}_{\text{index},t}} < \tau_g \,\right],
  \label{eq:gripper}
\end{equation}
with hysteresis on $\tau_g$ to suppress chatter at the threshold. This single
scalar is the entire hand representation beyond the wrist frame, following the
convention established for parallel-jaw targets in \cite{opentv, ace}.

\section{Object-Centric Representation}
\label{sec:objcentric}

Let $\bT^{W}_{O,t}$ denote the pose of the manipulated object in the workspace
frame. Each demonstration is stored in two forms. The workspace-anchored form
is $\bT^{W}_{H,t}$ directly; the object-relative form is
\begin{equation}
  \bT^{O}_{H,t} \;=\; \left(\bT^{W}_{O,t}\right)^{-1} \bT^{W}_{H,t}.
  \label{eq:objrel}
\end{equation}

The object-relative form is what transfers. If the same task is to be performed
with the object at a new pose $\bT^{R}_{O'}$ expressed in the robot frame, the
desired end-effector pose at time $t$ is
\begin{equation}
  \bT^{R,\mathrm{des}}_{E,t} \;=\;
  \bT^{R}_{O'} \; \bT^{O}_{H,t} \; \bT^{H}_{E},
  \label{eq:transfer}
\end{equation}
where $\bT^{H}_{E}$ is a fixed offset relating the human wrist frame to the
robot end-effector frame, determined once per embodiment. Trajectories are
retained densely rather than reduced to keyframes, since the continuous
velocity profile of the demonstration carries task-relevant information---
approach speed, contact timing---that sparse keyframes discard.

Both representations are exported. Storing both costs little and defers to the
downstream consumer a choice that the literature does not settle universally:
object-relative trajectories generalise across object placement
\cite{orion, spot} but depend on object pose estimation, whereas
workspace-anchored trajectories are simpler and adequate within a fixed layout.

\section{Retargeting as a Single Cost Functional}
\label{sec:retarget}

\subsection{Formulation}

Let $\bq_t \in \R^{n}$ be the manipulator configuration at time $t$ and
$f(\bq_t) \in \SE$ the forward kinematics of the end-effector. The task-space
residual with respect to the desired pose \eqref{eq:transfer} is the
logarithmic map of the pose error,
\begin{equation}
  \be_t(\bq_t) \;=\;
  \Log\!\left(
    \left(\bT^{R,\mathrm{des}}_{E,t}\right)^{-1} f(\bq_t)
  \right)^{\vee} \in \R^{6}.
  \label{eq:residual}
\end{equation}

Retargeting is then posed as the single constrained problem
\begin{equation}
\begin{aligned}
  \min_{\bq_{1:T}} \;\; J(\bq_{1:T}) \;=\; \sum_{t=1}^{T} \Big[\;
   & \underbrace{\omega_t \,\rho_\delta\!\left(
       \norm{\be_t(\bq_t)}_{\bm{\Sigma}}\right)}_{\text{data fidelity}}
   \;+\;
     \underbrace{\lambda_v \norm{\bq_t - \bq_{t-1}}^{2}}_{\text{velocity}}
   \\[1mm]
   &+\;
     \underbrace{\lambda_a \norm{\bq_t - 2\bq_{t-1} + \bq_{t-2}}^{2}}
       _{\text{acceleration}}
   \;+\;
     \underbrace{\lambda_r \norm{\bq_t - \bq_{\mathrm{ref}}}^{2}}
       _{\text{posture}}
  \;\Big]
\end{aligned}
\label{eq:cost}
\end{equation}
subject to
\begin{align}
  \bq_{\min} \;\le\; &\bq_t \;\le\; \bq_{\max}, \label{eq:cjoint}\\
  \norm{\bq_t - \bq_{t-1}} \;&\le\; \dot{\bq}_{\max}\,\Delta t,
    \label{eq:cvel}\\
  h_j(\bq_t) \;&\ge\; 0, \quad j = 1,\dots,m, \label{eq:ccoll}
\end{align}
where

\begin{center}
\begin{tabular}{@{\hspace{10mm}}p{28mm}@{}p{85mm}}
$\omega_t$ & per-frame perception confidence weight, defined in \eqref{eq:conf};\\
$\rho_\delta$ & Huber loss with transition parameter $\delta$;\\
$\bm{\Sigma}$ & diagonal matrix normalising translational and rotational residual units;\\
$\lambda_v, \lambda_a, \lambda_r$ & weights on velocity, acceleration and posture regularisation;\\
$h_j$ & collision and self-collision barrier functions.
\end{tabular}
\end{center}

\subsection{Confidence weighting}

The per-frame weight aggregates the landmark confidences that produced the
wrist estimate at time $t$,
\begin{equation}
  \omega_t \;=\;
  \left(\frac{1}{|\mathcal{I}|}\sum_{i \in \mathcal{I}}
        \max_k w^{k}_{i,t}\right)^{\alpha},
  \label{eq:conf}
\end{equation}
with $\mathcal{I}$ the landmark subset defining the wrist frame and $\alpha$
controlling how sharply low-confidence frames are discounted. The construction
is the offline analogue of the distance-dependent switching weights used in
online teleoperation retargeting \cite{opentv}: frames in which the hand is
partially occluded or poorly observed exert proportionally less influence on
the solution, and the smoothness terms carry the trajectory through them
instead of the data term dragging it toward a spurious observation.

\subsection{Robustification}

The Huber loss $\rho_\delta$ bounds the influence of gross outliers---a
mis-detected hand, a depth dropout that survives the confidence weighting---at
the level of the data term itself, rather than allowing a filtering stage to
smear the outlier across neighbouring frames \cite{dexteleop0}. This is the
mechanism that replaces pre-smoothing: outliers are down-weighted where they
enter, not attenuated after they have already propagated.

\subsection{Why a single functional}

Equation~\eqref{eq:cost} differs from the sequential alternative in one
structurally important respect. The smoothness penalties act on $\bq_t$, in the
space where smoothness is required, and they are traded against data fidelity
by the solver rather than applied blindly beforehand. In the sequential
pipeline, task-space filtering cannot guarantee joint-space smoothness because
the forward kinematics is nonlinear, and it cannot distinguish task-relevant
fast motion from noise because it operates without reference to the data term.
Here the trade-off is explicit: on high-confidence frames $\omega_t$ is large
and the solution tracks the observation; on low-confidence frames the
smoothness terms dominate and the solution interpolates.

Second-order smoothness, expressed by the acceleration term, follows the
$C^2$-continuity argument of DexFlow \cite{dexflow}. A trajectory that is
first-order smooth may still exhibit curvature discontinuities that appear as
jerk on the physical manipulator and as label noise to a policy trained on the
data.

\subsection{Solution}

The problem is solved with Pink on Pinocchio \cite{pink}. Each term of
\eqref{eq:cost} maps onto a weighted task in the quadratic program---a frame
task for the data term, posture and velocity regularisation tasks for the
smoothness terms---while \eqref{eq:cjoint}--\eqref{eq:ccoll} enter as
inequality constraints, with collision avoidance expressed through control
barrier functions. Levenberg--Marquardt damping stabilises the solution near
singularities.

Because the setting is offline, the problem is solved over the whole trajectory
rather than frame by frame, so the acceleration term is well defined at every
interior timestep and the solution is not constrained by the causality
requirement that an online solver faces. The trajectory is initialised from the
previous solution warm-started along $t$, and boundary conditions at $t=1,2$
are handled by replicating the initial configuration.

\subsection{Interpolation}

Where gaps remain after fusion---landmarks lost across all views---the
trajectory is completed by cubic spline interpolation over $\SE$ prior to
optimisation. Dynamic filters that impose a motion model, such as Kalman or
extended Kalman filters, are deliberately not applied: they introduce a prior
about the dynamics that is subsequently re-estimated by the optimiser, and the
interaction between the two priors is difficult to characterise. The optimiser
already imposes the required regularity through
$\lambda_v$ and $\lambda_a$.

\section{Replay Validation}
\label{sec:replay}

Synthesised trajectories are executed on the physical manipulator before the
dataset is finalised. The stage serves three purposes.

First, it filters infeasibility. A trajectory may satisfy
\eqref{eq:cjoint}--\eqref{eq:ccoll} in the kinematic model and still fail on
hardware through controller tracking limits, unmodelled collisions or workspace
boundary conditions. Such episodes are rejected rather than exported.

Second, it closes the embodiment gap at the source. The proprioceptive states
recorded during replay are those the robot actually attained, not those the
optimiser predicted. Exporting the recorded rather than the synthesised states
means the dataset contains honest observation--action pairs from the target
embodiment, which is precisely the property that teleoperated datasets possess
by construction and that naive video-derived datasets lack.

Third, it provides a curation signal. Episodes are ranked by tracking residual
during replay, and the ranking is available to downstream consumers as a
quality score, in the spirit of rollout-based demonstration curation
\cite{demoscore} but applied before rather than after policy training.

Failure of an episode triggers either rejection or re-solution of
\eqref{eq:cost} with modified weights---typically an increased $\lambda_v$ or a
relaxed data weight in the region of failure---as indicated by the feedback
edge in Fig.~\ref{fig:pipeline}. The number of re-solution attempts is bounded
to prevent overfitting the trajectory to the replay controller.

\section{Dataset Export}
\label{sec:export}

Surviving episodes are written in a LeRobot-compatible HDF5 layout. Each
episode contains the synchronised camera streams, the recorded joint states and
gripper states from replay, the object-relative and workspace-anchored wrist
trajectories, per-frame confidence weights, and the replay tracking residual.
Retaining the confidence weights and residuals in the exported artefact allows
downstream curation to be repeated with different thresholds without
re-recording.

\section{Assumptions and Simplifications}
\label{sec:assumptions}

The following assumptions delimit the scope of the work and are revisited in
Section~\ref{sec:limitations}.

\begin{enumerate}[leftmargin=2em]
  \item \emph{Parallel-jaw target.} The hand is reduced to a wrist pose and a
        binary gripper state. Multi-fingered dexterous targets would require
        the finger-level representation rejected in
        Section~\ref{sec:lit-hand}.
  \item \emph{Rigid objects.} The object-relative representation
        \eqref{eq:objrel} presupposes a rigid body with a well-defined pose.
        Deformable objects would require a flow-based representation
        \cite{novaflow}.
  \item \emph{Single-arm, quasi-static tasks.} The task class is restricted to
        single-arm tabletop manipulation at moderate speed, where the
        quasi-static assumption underlying purely kinematic retargeting holds.
  \item \emph{Fixed workspace.} The ChArUco anchor is rigidly attached to the
        workspace and the camera rig is not moved within a recording session.
  \item \emph{Kinematic-only transfer.} Forces and contact dynamics are not
        estimated. Tasks whose success depends on force regulation are outside
        the scope.
\end{enumerate}
